\documentclass[conference]{IEEEtran}
\IEEEoverridecommandlockouts
% The preceding line is only needed to identify funding in the first footnote. If that is unneeded, please comment it out.
\usepackage{cite}
\usepackage{amsmath,amssymb,amsfonts}
\usepackage{algorithmic}
\usepackage{graphicx}
\usepackage{textcomp}
\usepackage{xcolor}
\def\BibTeX{{\rm B\kern-.05em{\sc i\kern-.025em b}\kern-.08em
    T\kern-.1667em\lower.7ex\hbox{E}\kern-.125emX}}

\begin{document}

\title{Klasifikasi Multi-Label untuk Pengenalan Huruf Hijaiyah Sambung Menggunakan Transfer Learning\\
{\footnotesize \textsuperscript{*}Proposal Skripsi Program Studi S-1 Informatika}
}

\author{\IEEEauthorblockN{Muchamad Budi Fabiantoro}
\IEEEauthorblockA{\textit{Program Studi Informatika, Fakultas Teknologi dan Industri} \\
\textit{Universitas Internasional Semen Indonesia (UISI)}\\
Gresik, Indonesia \\
NIM: 22314050111003}
\and
\IEEEauthorblockN{Moch. Zain Hasan, S.Kom., M.Kom.}
\IEEEauthorblockA{\textit{Departemen Informatika, Fakultas Teknologi dan Industri} \\
\textit{Universitas Internasional Semen Indonesia (UISI)}\\
Gresik, Indonesia \\
Dosen Pembimbing}
}

\maketitle

\begin{abstract}
Penguasaan penulisan huruf Hijaiyah bersambung pada kurikulum Iqro 3 merupakan fondasi utama dalam pembelajaran literasi Al-Qur'an tingkat dasar. Selama ini, evaluasi kemampuan menulis huruf Hijaiyah masih bergantung pada koreksi visual manual oleh guru secara tatap muka, yang menghadirkan kendala keterbatasan akses bimbingan mandiri dan konsistensi penilaian. Penelitian ini mengusulkan pengembangan model pengenalan tulisan tangan huruf Hijaiyah bersambung secara nirsegmentasi (\textit{segmentation-free}) menggunakan arsitektur gabungan \textit{Convolutional Neural Network} dan \textit{Long Short-Term Memory} (CNN-LSTM) dengan strategi \textit{transfer learning} berbasis \textit{backbone} MobileNetV3. Masukan sistem berupa citra digital tunggal hasil ekspor \textit{canvas digital} dengan batas kosakata kata sambung pendek (2--3 karakter) mengacu pada buku Iqro 3. Citra masukan diekstraksi representasi ciri spasialnya oleh MobileNetV3, kemudian dimodelkan urutan temporalnya dari arah kanan ke kiri oleh LSTM, dan didekodekan menggunakan fungsi \textit{Connectionist Temporal Classification} (CTC). Kinerja model dievaluasi secara komprehensif menggunakan metrik \textit{exact sequence match accuracy}, akurasi per posisi karakter, dan \textit{Character Error Rate} (CER). Pendekatan ini diharapkan mampu menghadirkan sistem kecerdasan buatan yang ringan, akurat, dan efisien sebagai teknologi pembantu (\textit{assisting technology}) untuk evaluasi mandiri dalam pembelajaran literasi Al-Qur'an.
\end{abstract}

\begin{IEEEkeywords}
Hijaiyah Sambung, CNN-LSTM, Transfer Learning, MobileNetV3, Segmentation-Free, Canvas Digital, Iqro 3, OCR Luring.
\end{IEEEkeywords}

\section{Pendahuluan}

\subsection{Latar Belakang}
Di era transformasi digital dan perkembangan kecerdasan buatan (\textit{Artificial Intelligence}) saat ini, teknologi informasi berkembang sangat pesat menjadi instrumen pembantu (\textit{assisting technology}) yang mempermudah berbagai aktivitas manusia, termasuk dalam bidang pendidikan keagamaan. Integrasi teknologi cerdas kini diarahkan untuk menghadirkan media pembelajaran yang adaptif, interaktif, dan dapat diakses secara mandiri oleh peserta didik kapan pun dan di mana pun. Di Indonesia, pembelajaran membaca dan menulis aksara Al-Qur'an pada tingkat dasar umumnya diselenggarakan menggunakan metode Iqro. Pada tahap awal (buku Iqro 1 dan Iqro 2), peserta didik mempelajari huruf-huruf Hijaiyah dalam bentuk tunggal yang terisolasi (\textit{isolated characters}). Ketika melangkah ke buku \textbf{Iqro 3}, peserta didik mulai diperkenalkan pada kata-kata pendek yang tersusun atas dua hingga tiga huruf Hijaiyah yang ditulis secara bersambung. Penguasaan penulisan huruf Hijaiyah bersambung memegang peranan yang sangat fundamental karena seluruh teks Al-Qur'an dan literatur berbahasa Arab disajikan dalam bentuk rangkaian huruf kursif bersambung.

Meskipun memegang peranan krusial, proses evaluasi kemampuan menulis huruf Hijaiyah sejauh ini masih sangat bergantung pada koreksi visual manual oleh guru atau ustaz secara tatap muka langsung di kelas. Sebagaimana digarisbawahi oleh Rahmatulloh dkk. \cite{rahmatulloh2022}, proses pembelajaran huruf Hijaiyah konvensional menuntut kehadiran pembimbing khusus secara langsung untuk memverifikasi kebenaran setiap goresan huruf. Ketergantungan mutlak pada evaluasi manual ini menimbulkan kendala nyata, terutama ketika kegiatan pembelajaran berlangsung secara jarak jauh (\textit{remote learning}), membatasi frekuensi latihan mandiri peserta didik di luar jam pertemuan kelas, serta menghadirkan keterbatasan objektivitas dan konsistensi penilaian pada skala kelas yang berukuran besar. Mengingat laju perkembangan zaman yang serba cepat dan menuntut solusi praktis, Rahmatulloh dkk. \cite{rahmatulloh2022} menegaskan urgensi kehadiran teknologi cerdas berupa sistem pengenalan tulisan tangan (\textit{Optical Character Recognition} / OCR) luring berbasis \textit{deep learning} yang mampu mengidentifikasi tulisan dan memandu proses pembelajaran secara otomatis dengan memberikan umpan balik koreksi seketika (\textit{instant warning and feedback}) apabila terjadi kesalahan penulisan.

Peluang penerapan teknologi pembantu evaluasi mandiri ini semakin terbuka luas seiring meluasnya adopsi perangkat layar sentuh (\textit{touchscreen}) seperti tablet dan ponsel pintar yang dilengkapi fasilitas media \textit{canvas digital} interaktif. Melalui media tersebut, peserta didik dapat berlatih menuliskan aksara Hijaiyah langsung di atas layar gawai dengan panduan garis titik-titik (\textit{tracing dots}), di mana hasil goresan tersebut dapat diekspor menjadi berkas citra digital 2 dimensi (2D). Pemanfaatan perangkat layar sentuh untuk latihan menulis karakter Arab telah terbukti diterima secara antusias dan efektif diterapkan pada anak-anak usia sekolah dasar maupun kalangan pendidik, sebagaimana ditunjukkan pada evaluasi aplikasi tablet \textit{Mutqin} oleh Durayhim dkk. \cite{durayhim2023}. Citra hasil goresan pada \textit{canvas digital} tersebut memiliki potensi yang sangat besar untuk dijadikan masukan (\textit{input}) bagi sistem kecerdasan buatan guna menyediakan evaluasi otomatis yang cepat, konsisten, dan terukur.

Sebagai upaya mewujudkan sistem pengenalan tulisan tangan otomatis, penelitian acuan utama oleh Rahmatulloh dkk. \cite{rahmatulloh2022} yang berjudul \textit{``Optimization of Hijaiyah Letter Handwriting Recognition Model Based on Deep Learning''} telah membuktikan efektivitas pendekatan \textit{Deep Learning}, khususnya arsitektur \textit{Convolutional Neural Network} (CNN) yang dioptimalkan dengan algoritma ADAM, dalam mengenali dan mengklasifikasikan tulisan tangan huruf Hijaiyah. Penelitian tersebut sukses mencatatkan tingkat akurasi klasifikasi yang tinggi, yakni sebesar 91{,}24\% pada dataset tulisan tangan anak-anak usia 7--12 tahun (Hijja2) dan 98{,}42\% pada dataset tulisan dewasa (AHCD). Di akhir penelitiannya, Rahmatulloh dkk. \cite{rahmatulloh2022} memberikan arahan dan rekomendasi penting bagi penelitian selanjutnya (\textit{future work}), yaitu mengembangkan dan mengimplementasikan model pengenalan tulisan tangan Hijaiyah tersebut ke dalam aplikasi berbasis perangkat bergerak (\textit{mobile}) agar dapat digunakan secara praktis untuk membantu proses pembelajaran secara interaktif.

Kendati demikian, model yang dibangun oleh Rahmatulloh dkk. \cite{rahmatulloh2022}---serta sebagian besar penelitian pengenalan aksara Arab di Indonesia \cite{masruroh2023}---masih terbatas pada pengenalan huruf tunggal yang berdiri sendiri (\textit{isolated characters}). Di sisi lain, pada jenjang pembelajaran nyata kurikulum Iqro 3, peserta didik diwajibkan menuliskan rangkaian huruf yang saling bersambung membentuk kata. Pengenalan tulisan tangan huruf Hijaiyah bersambung menghadirkan tantangan komputasi yang jauh lebih rumit dibandingkan pengenalan huruf tunggal. Pertama, morfologi huruf Arab bersifat dinamis dan fleksibel; setiap karakter dapat berubah menjadi hingga empat bentuk glif visual yang berbeda bergantung pada posisinya di dalam kata, yaitu pada posisi awal (\textit{initial}), tengah (\textit{medial}), akhir (\textit{final}), atau berdiri sendiri (\textit{isolated}) \cite{altwaijry2021, durayhim2023}. Kedua, sebagaimana diidentifikasi oleh Rahmatulloh dkk. \cite{rahmatulloh2022}, banyak huruf Hijaiyah berbagi kerangka dasar yang identik dan hanya dibedakan oleh jumlah serta letak titik diakritik (misalnya kelompok \textit{Ba-Ta-Tsa} dan \textit{Jim-Ha-Kha}), di mana pada penulisan bersambung titik-titik ini rentan terdistorsi, bergeser, atau bertumpang tindih dengan goresan huruf di sekitarnya, padahal keberadaan titik sangat krusial karena dapat mengubah total identitas huruf dan makna kata \cite{albattah2022}. Ketiga, penulisan huruf yang saling bertaut membentuk sambungan alami (\textit{ligature}) tanpa batas pemisah yang jelas. Karakteristik ligatur ini menyebabkan pendekatan konvensional berbasis segmentasi karakter (\textit{segmentation-based}) kerap mengalami kegagalan fatal akibat pemotongan berlebih (\textit{over-segmentation}) atau kesalahan potong pada persilangan garis ligatur, di mana galat segmentasi tersebut merambat langsung menjadi kesalahan fatal dalam pengenalan karakter (\textit{cascading error}) \cite{peng2024, mosbah2024}. Kompleksitas ini semakin tinggi pada tulisan tangan anak-anak atau pemula yang memiliki variabilitas goresan tinggi, ketidakstabilan tekanan garis, dan deformasi proporsi bentuk yang signifikan, terbukti dari akurasi pengenalan dataset anak yang secara empiris konsisten lebih rendah (82{,}5\%--91\%) dibandingkan dataset tulisan orang dewasa (97\%--98\%) \cite{altwaijry2021, durayhim2023, alwagdani2023, masruroh2023}.

Guna mengatasi kelemahan mendasar dari segmentasi fisik tersebut, paradigma pengenalan sekuensial nirsegmentasi (\textit{segmentation-free sequence recognition}) menjadi solusi yang paling mutakhir dan menjanjikan. Pendekatan ini memperlakukan citra kata sebagai barisan representasi visual utuh yang dipetakan secara langsung menjadi urutan teks (\textit{sequence of characters}). Integrasi arsitektur hibrida antara \textit{Convolutional Neural Network} (CNN) dan \textit{Long Short-Term Memory} (LSTM) terbukti sangat tangguh dalam mengenali barisan teks tulisan tangan bersambung \cite{saini2024, peng2024}. Pada skema hibrida ini, lapisan CNN bertindak mengekstraksi representasi peta ciri visual spasial dari citra masukan, sedangkan lapisan rekuren LSTM membaca dependensi kontekstual urutan karakter secara temporal dari arah kanan ke kiri sesuai arah alami penulisan aksara Arab, sehingga mampu menyelesaikan persoalan ligatur yang rumit tanpa memerlukan pemotongan karakter secara fisik \cite{mosbah2024}.

Selain tantangan morfologi sambung, kendala teknis penting dalam merealisasikan rekomendasi Rahmatulloh dkk. \cite{rahmatulloh2022} menuju aplikasi interaktif adalah keterbatasan jumlah data latih citra \textit{canvas digital} mandiri dan tuntutan efisiensi komputasi pada perangkat bergerak. Melatih jaringan konvolusi dalam (\textit{deep CNN}) dari awal (\textit{from scratch}) dengan data berskala terbatas sangat rentan menimbulkan \textit{overfitting} serta memperlambat konvergensi model. Dalam konteks ini, penerapan \textit{transfer learning} menjadi strategi krusial untuk memanfaatkan bobot yang telah dipelajari oleh model terlatih sebelumnya sebagai pengekstraksi ciri visual yang kaya dan tergeneralisasi \cite{lamtougui2024, masruroh2023}. Sejalan dengan tren teknologi masa kini yang menekankan komputasi cerdas pada perangkat berdaya rendah (\textit{mobile and edge computing}), pemilihan arsitektur \textit{backbone} modern yang sangat efisien seperti \textbf{MobileNetV3} (\textit{Large}) menawarkan keunggulan latensi sangat rendah, bobot parameter yang ringkas (sekitar 5,4 juta parameter), dan akurasi tinggi berkat pemanfaatan blok \textit{inverted residual} dengan modul \textit{Squeeze-and-Excitation} (SE) serta aktivasi \textit{hard-swish} \cite{howard2019, elkhayati2025}. Karakteristik ini sangat ideal untuk mendukung inferensi cepat secara luring pada perangkat pembelajaran berbasis \textit{canvas digital}.

Kajian literatur menunjukkan bahwa penelitian yang mengkaji tulisan tangan Arab bersambung sejauh ini umumnya masih berfokus pada naskah dokumen pindaian kertas statis tulisan orang dewasa seperti dataset IFN/ENIT dengan pendekatan klasifikasi kata utuh (\textit{whole-word classification}) \cite{lamtougui2024, ahmed2021}. Berdasarkan telaah terhadap penelitian terdahulu, terdapat tiga kesenjangan penelitian (\textit{research gap}) utama yang belum terjawab:
\begin{enumerate}
    \item \textbf{Kesenjangan media masukan dan kurikulum:} Belum ada penelitian yang membangun dan mengevaluasi model pada citra hasil ekspor \textit{canvas digital} interaktif yang secara terstruktur dirancang berdasarkan kosakata kurikulum pembelajaran Al-Qur'an tingkat dasar (\textbf{Iqro 3}).
    \item \textbf{Kesenjangan dekomposisi sekuensial:} Penelitian tulisan tangan Arab bersambung sebelumnya \cite{lamtougui2024, ahmed2021} memperlakukan kata sebagai satu label kelas utuh statis (\textit{whole-word}), bukan memetakan citra menjadi urutan label karakter individual $[H_1, H_2, H_3]$ secara sekuensial nirsegmentasi sehingga model dapat mengenali kombinasi huruf baru secara fleksibel.
    \item \textbf{Kesenjangan integrasi metode:} Penelitian yang memanfaatkan \textit{transfer learning} \cite{masruroh2023, durayhim2023, lamtougui2024} terbatas pada CNN murni tanpa pemodel urutan LSTM, sedangkan penelitian CNN-LSTM \cite{saini2024, peng2024} hanya menguji aksara Latin atau pindaian dokumen cetak. Belum ada penelitian yang menggabungkan \textit{transfer learning} MobileNetV3 sebagai pengekstraksi ciri spasial ringan dengan LSTM dan dekoder CTC (\textit{Connectionist Temporal Classification}) untuk pengenalan kata sambung pendek huruf Hijaiyah.
\end{enumerate}

Berdasarkan latar belakang permasalahan, arah perkembangan teknologi pembantu (\textit{assisting technology}), serta pijakan riset acuan Rahmatulloh dkk. \cite{rahmatulloh2022}, penelitian ini mengusulkan pengembangan model pengenalan sekuensial berbasis arsitektur gabungan \textbf{CNN-LSTM dengan \textit{transfer learning} MobileNetV3} untuk mengenali tulisan tangan huruf Hijaiyah bersambung pada materi Iqro 3. Model ini dirancang untuk memproses citra digital hasil ekspor \textit{canvas digital} secara \textit{segmentation-free} dan memprediksi urutan karakter dari kanan ke kiri secara tepat. Penelitian ini difokuskan pada perancangan dataset citra dan rekayasa model \textit{machine learning}, dengan harapan dapat menghadirkan sistem kecerdasan buatan yang andal, akurat, dan efisien dalam mendukung evaluasi mandiri literasi Al-Qur'an tingkat dasar.

\subsection{Rumusan Masalah}
Berdasarkan latar belakang permasalahan yang telah diuraikan, rumusan masalah dalam penelitian ini dirumuskan sebagai berikut:
\begin{enumerate}
    \item Bagaimana merancang dan membangun dataset citra tulisan tangan huruf Hijaiyah bersambung berbasis ekspor \textit{canvas digital} yang merepresentasikan seluruh kosakata kata sambung pada buku Iqro 3?
    \item Bagaimana mengimplementasikan arsitektur CNN-LSTM dengan memanfaatkan \textit{transfer learning} MobileNetV3 untuk mengenali urutan huruf Hijaiyah yang ditulis bersambung pada citra tanpa melalui tahapan segmentasi manual (\textit{segmentation-free})?
    \item Bagaimana kinerja model CNN-LSTM dengan \textit{transfer learning} MobileNetV3 dalam memprediksi urutan karakter huruf Hijaiyah bersambung berdasarkan metrik evaluasi pengenalan sekuensial (\textit{exact sequence match accuracy}, akurasi per karakter, dan \textit{Character Error Rate})?
\end{enumerate}

\subsection{Tujuan Penelitian}
Sesuai dengan rumusan masalah yang diajukan, tujuan yang ingin dicapai dalam pelaksanaan penelitian ini adalah:
\begin{enumerate}
    \item Merancang dan membangun dataset citra tulisan tangan huruf Hijaiyah bersambung berbasis ekspor \textit{canvas digital} yang terstruktur dan representatif mengacu pada daftar kosakata buku Iqro 3.
    \item Mengembangkan dan mengimplementasikan model pengenalan tulisan tangan bersambung berbasis arsitektur gabungan CNN-LSTM dengan \textit{transfer learning} MobileNetV3 yang mampu memetakan citra masukan secara langsung menjadi urutan karakter dari arah kanan ke kiri secara nirsegmentasi (\textit{segmentation-free}).
    \item Menguji, menganalisis, dan mengevaluasi performa model CNN-LSTM dengan \textit{transfer learning} MobileNetV3 dalam mengenali urutan karakter huruf Hijaiyah bersambung menggunakan metrik \textit{exact sequence match accuracy}, akurasi per karakter, dan \textit{Character Error Rate} (CER).
\end{enumerate}

\subsection{Manfaat Penelitian}
Hasil penelitian ini diharapkan dapat memberikan kontribusi nyata yang terbagi ke dalam:
\subsubsection{Manfaat Teoretis}
Memberikan kontribusi ilmiah dalam pengembangan metode pengenalan teks sekuensial nirsegmentasi berbasis hibrida CNN-LSTM dengan strategi \textit{transfer learning} pada domain aksara kursif non-Latin, serta memperkaya literatur pemrosesan citra terkait variasi morfologi posisional dan ligatur kata pendek.

\subsubsection{Manfaat Praktis}
\begin{itemize}
    \item \textbf{Bagi Peserta Didik:} Menyediakan asisten pembantu evaluasi penulisan kata Hijaiyah secara mandiri, objektif, dan instan di luar kelas.
    \item \textbf{Bagi Pendidik:} Membantu proses asesmen tulisan tangan secara konsisten serta meringankan beban pemeriksaan visual manual pada skala kelas besar.
    \item \textbf{Bagi Pengembang:} Menyediakan arsitektur model \textit{deep learning} yang ringan (MobileNetV3) dan siap diintegrasikan ke modul evaluasi aplikasi bergerak (\textit{mobile/tablet}).
\end{itemize}

\subsubsection{Manfaat Akademis}
Menjadi rujukan ilmiah dan komparasi empiris bagi peneliti lanjutan dalam domain OCR aksara Arab/Hijaiyah di Indonesia, khususnya kelanjutan dari studi karakter terisolasi menuju teks bersambung.

\subsection{Batasan Penelitian}
Agar ruang lingkup penelitian tetap terarah dan terukur, batasan penelitian ditetapkan sebagai berikut:
\begin{enumerate}
    \item \textbf{Objek Penelitian:} Citra tulisan tangan huruf Hijaiyah bersambung 2D berformat PNG/JPEG hasil ekspor \textit{canvas digital} (pengenalan luring/\textit{offline}), bukan data runtun waktu koordinat goresan (\textit{online stroke trajectory}).
    \item \textbf{Data Penelitian:} Diperoleh dari latihan penulisan mandiri pada \textit{canvas digital} dengan alat bantu pola titik-titik panduan (\textit{tracing dots}); kosakata dibatasi pada buku Iqro 3 (panjang 2--3 huruf per citra); \textit{whitelist} 28 huruf Hijaiyah dasar tanpa harakat atau tanda baca; pembagian data menggunakan skema \textit{participant-wise split}.
    \item \textbf{Metode dan Pendekatan:} Pendekatan nirsegmentasi berbasis CNN-LSTM; komponen CNN menggunakan MobileNetV3 (\textit{Large}) terbobot ImageNet (\textit{transfer learning}); LSTM searah untuk pemodelan urutan kanan ke kiri; dekoder CTC untuk inferensi urutan teks.
    \item \textbf{Variabel dan Parameter:} Input citra greyscale $256 \times 64$ piksel; output barisan string 2--3 karakter; optimasi menggunakan algoritma ADAM.
    \item \textbf{Proses dan Tahapan:} Pengumpulan data berbasis \textit{prompt}, kendali mutu data (QC), pra-pemrosesan citra, pelatihan dua tahap (\textit{two-stage training}), dan pengujian model.
    \item \textbf{Metrik Evaluasi:} Diukur menggunakan \textit{Exact Sequence Match Accuracy}, \textit{Character/Positional Accuracy}, dan \textit{Character Error Rate} (CER).
    \item \textbf{Aspek di Luar Cakupan:} Mahasiswa berfokus pada dataset dan model (gamifikasi dan integrasi aplikasi dilakukan oleh kolaborator); tidak mengevaluasi kaidah hukum tajwid, panjang-pendek bacaan (\textit{mad}), maupun estetika kaligrafi (\textit{khat}).
\end{enumerate}

\subsection{Sistematika Penulisan}
Proposal penelitian ini disusun dalam tiga bab utama:
\begin{itemize}
    \item \textbf{Bab I Pendahuluan:} Memuat latar belakang, perumusan masalah, tujuan, manfaat, batasan penelitian, dan sistematika penulisan.
    \item \textbf{Bab II Kajian Pustaka dan Landasan Teori:} Membahas telaah penelitian terdahulu, analisis \textit{research gap}, teori morfologi Hijaiyah bersambung, MobileNetV3, LSTM, dan mekanisme nirsegmentasi CTC.
    \item \textbf{Bab III Metodologi Penelitian:} Menguraikan alur penelitian, protokol akuisisi dataset \textit{canvas}, tahap prapemrosesan citra, rancangan arsitektur CNN-LSTM MobileNetV3, skenario pelatihan dua tahap, dan rencana evaluasi.
\end{itemize}

\begin{thebibliography}{00}

\bibitem{rahmatulloh2022}
A.~Rahmatulloh, R.~Rizal, R.~I. Gunawan, I.~Darmawan, and B.~Z. Rahmat, ``Optimization of Hijaiyah letter handwriting recognition model based on deep learning,'' in \emph{Proc. 2022 Int. Conf. Advancement in Data Science, E-learning and Information Systems (ICADEIS)}, 2022, pp. 1--6, doi: 10.1109/ICADEIS56544.2022.10037496.

\bibitem{durayhim2023}
A.~B. Durayhim, A.~Al-Ajlan, I.~Al-Turaiki, and N.~Altwaijry, ``Towards accurate children's Arabic handwriting recognition via deep learning,'' \emph{Appl. Sci.}, vol. 13, no. 3, p. 1692, 2023, doi: 10.3390/app13031692.

\bibitem{masruroh2023}
S.~U. Masruroh, M.~F. Syahid, F.~Munthaha, A.~T. Muharram, and R.~A. Putri, ``Deep convolutional neural networks transfer learning comparison on Arabic handwriting recognition system,'' \emph{Int. J. Informatics Visualiz.}, vol. 7, no. 2, pp. 330--337, 2023.

\bibitem{altwaijry2021}
N.~Altwaijry and I.~Al-Turaiki, ``Arabic handwriting recognition system using convolutional neural network,'' \emph{Neural Comput. Appl.}, vol. 33, no. 7, pp. 2249--2261, 2021, doi: 10.1007/s00521-020-05070-8.

\bibitem{albattah2022}
W.~Albattah and S.~Albahli, ``Intelligent Arabic handwriting recognition using different standalone and hybrid CNN architectures,'' \emph{Appl. Sci.}, vol. 12, no. 19, p. 10155, 2022, doi: 10.3390/app121910155.

\bibitem{alwagdani2023}
M.~S. Alwagdani and E.~S. Jaha, ``Deep learning-based child handwritten Arabic character recognition and handwriting discrimination,'' \emph{Sensors}, vol. 23, no. 15, p. 6774, 2023, doi: 10.3390/s23156774.

\bibitem{peng2024}
G.~Peng, ``Segmentation-free recognition algorithm based on deep learning for handwritten text image,'' \emph{J. Adv. Inf. Technol.}, vol. 4, no. 2, pp. 169--178, 2024, doi: 10.37965/jait.2024.0473.

\bibitem{mosbah2024}
L.~Mosbah, I.~Moalla, T.~M. Hamdani, B.~Neji, T.~Beyrouthy, and A.~M. Alimi, ``ADOCRNet: A deep learning OCR for Arabic documents recognition,'' \emph{IEEE Access}, vol. 12, pp. 55620--55629, 2024, doi: 10.1109/ACCESS.2024.3379530.

\bibitem{saini2024}
M.~L. Saini, D.~C. Sati, R.~S. Telikicharla, and Mahadev, ``Handwritten English script recognition system using CNN and LSTM,'' in \emph{Proc. 2024 IEEE Int. Conf. Contemporary Computing and Communications (InC4)}, 2024, pp. 1--6, doi: 10.1109/InC460750.2024.10649099.

\bibitem{lamtougui2024}
H.~Lamtougui, H.~El~Moubtahij, H.~Fouadi, and K.~Satori, ``Improving Arabic handwritten text recognition through transfer learning with convolutional neural network-based models,'' \emph{Bull. Electr. Eng. Informatics}, vol. 13, no. 6, pp. 4294--4305, 2024, doi: 10.11591/eei.v13i6.8178.

\bibitem{howard2019}
A.~Howard \emph{et~al.}, ``Searching for MobileNetV3,'' in \emph{Proc. IEEE/CVF Int. Conf. Computer Vision (ICCV)}, 2019, pp. 1314--1324, doi: 10.1109/ICCV.2019.00140.

\bibitem{elkhayati2025}
M.~El~Khayati, A.~Maafiri, Y.~Himeur, H.~A. Alkhazaleh, S.~Atalla, and W.~Mansoor, ``Leveraging transfer learning and mobile-enabled convolutional neural networks for improved Arabic handwritten character recognition,'' \emph{arXiv:2509.05019}, 2025.

\bibitem{ahmed2021}
R.~Ahmed \emph{et~al.}, ``Novel deep convolutional neural network-based contextual recognition of Arabic handwritten scripts,'' \emph{Entropy}, vol. 23, no. 3, p. 340, 2021, doi: 10.3390/e23030340.

\end{thebibliography}

\end{document}
